\section{Conclusion}\label{sec:conclusion}

A decade of empirical work has established that restricting a procurement auction to small firms raises winning prices. This paper goes a step further: it identifies which mechanism produces the markup under explicit modeling assumptions and prices the welfare consequence against a specific policy alternative. The \policyCutoffMonth{} extension of S\~ao Paulo's SME-only regime to medical supplies supplies the identifying variation. The structural framework---asymmetric IPV with Krasnokutskaya UH cleaning, Haile-Tamer drop-out point-ID, and observed-equilibrium-entry treatment in the Athey-Seira spirit---is established in the literature; the paper's contribution is to combine these conditions with a single clean legal trigger that the partial-preference structural literature has not had access to. The exercise runs under observed equilibrium entry, with strict primitive invariance reported as a bound rather than the main specification.

Five points organize the contribution. First, the SME-only set-aside imposes an annual welfare loss on the order of R\$\,\welfAnnualBRLLo--\welfAnnualBRLHi~million (US\$\,\welfAnnualUSDLo--\welfAnnualUSDHi~million) on a single product group of S\~ao Paulo's R\$\,\becAnnualVolBRL~billion centralized platform; the welfare cost (allocative DWL plus MCPF distortion at $\lambda = \lambdaMain$) reaches \welfLossPctLthirtyNp~percent of $p_{S_1}$ in non-pharma and \welfLossPctLthirtyPh~percent in pharma. Second, a \optPrefThreshold{} SME price preference welfare-dominates the set-aside in thick standardized non-pharmaceutical markets across both the main and the strict-invariance specifications, at near-zero fiscal cost. Third, endogenous SME entry is the government's partial offset rather than the source of the markup: a fixed-pool counterfactual would deliver a simulated price effect roughly \latentShockRatioNp~percent larger than the realized one in non-pharmaceuticals and \latentShockRatioPh~percent larger in pharmaceuticals. Fourth, the simulated set-aside price effect operates predominantly within the auction under the model's main specification: the within-auction component---which I label \emph{sheltered bidding}---accounts for roughly three-quarters of the simulated total effect, in a window stable across the cost-distribution estimators reported. Fifth, in thin heterogeneous pharmaceutical markets the ranking against the price preference depends on how the post-policy SME pool is modeled; the conditional reading is itself a substantive result about where bidder exclusion delivers proportional returns and where it does not.

The policy reading follows. The question is not whether bidder exclusion is universally good or bad, but when: a moderate price preference is recommended where the protected pool is thick enough and the good standardized enough that it recovers nearly the same distributive gain at a fraction of the welfare cost; bidder exclusion remains defensible only where the pool is thin, the good heterogeneous, and the planner's welfare weight on SME surplus is high enough to absorb the larger allocative wedge. The choice depends on equilibrium-selection details this paper identifies but does not resolve. A general defense of bidder-exclusion rules across procurement categories is harder to sustain than current Brazilian and international practice would suggest; a uniform replacement of exclusion by preference is stronger than the data support. The multi-platform extension flagged in Section~\ref{sec:discussion} is the natural test of generalizability.

The structural exercise leaves three margins outside its scope---structural identification of entry costs, dynamic capacity accumulation, and the FPSB-with-preference equilibrium---each flagged as future work in Section~\ref{sec:discussion}. The choice belongs at the level of the product class, not the level of the procurement code.
